\section{Calculer avec des puissances}
\subsection{Définition du carré et du cube}
\begin{definition}[\it{Carré d'un nombre}]
    Le \bb{carré} d'un nombre $a$ est le produit de ce nombre par lui-même. 
    On le note $a^2$ et on lit  "\bb{a au carré}" ou \bb{"a exposant 2"} :
    $$a^2=a\times a$$
    Le petit 2 s'appelle l'\bb{exposant} : il indique le nombre de facteurs égaux à
    $a$ (\it{le nombre de fois où on multiplie $a$ par lui même})
\end{definition}

\begin{exe}
    \begin{itemize}
        \item $7^2\answer{ = 7\times 7 = 49}$. On dit : "\it{Le carré de 7 est 49}"
            \bb{ou} "\it{7 au carré vaut 49}"
        \item $12^2=\answer{12 \times 12 = 144}$
        \item $1.5^2=\answer{1.5 \times 1.5 = 2.25}$
    \end{itemize}
\end{exe}

\begin{exo}
    \begin{itemize}
        \item $9^2=\answer{9\times 9=81}$
        \item $12^2 = \answer{12\times 12=144}$
        \item $0.5^2 = \answer{ 0,5\times 0,5=0,25}$  
        \item $15^2 =  \answer{15\times 15=225}$
        \item $1,2^2 = \answer{1.2\times 1.2 = 1.44}$

    \end{itemize}
\end{exo}


\begin{attention}
    $6^2$ n'est pas $6\times 2$. On ne multiplie pas par l'exposant ! Il indique
    le nombre de facteur. $6^2=36$ alors que $6\times2=12$
\end{attention}

\begin{definition}[\it{Cube d'un nombre}]
    Le \bb{cube} d'un nombre $a$ est le produit de trois facteurs égaux à $a$. 
    On le note $a^3$ et on lit  "a au cube" ou "a exposant 3" :
    $$a^3=a\times a\times a$$
\end{definition}

\begin{exe}
    \begin{itemize}
        \item $7^2\answer{ = 7\times 7 = 49}$. On dit : "\it{Le carré de 7 est 49}"
            \bb{ou} "\it{7 au carré vaut 49}"
        \item $12^2=\answer{12 \times 12 = 144}$
        \item $1.5^2=\answer{1.5 \times 1.5 = 2.25}$
    \end{itemize}
\end{exe}

\begin{exo}
    \begin{itemize}
        \item $2^3=\answer{\underline{2 \times 2} \times 2 = 4 \times 2 = 8}$
        \item $4^3=\answer{\underline{4\times 4}\times 4 = 16\times 4 = 64}$
        \item $10^3=\answer{\underline{10\times 10}\times 10 = 100\times 10 = 1000}$
        \item $0.1^3=\answer{\underline{0.1\times 0.1}\times 0.1 = 0.01\times 0.1= 0.001}$
        \item $5^3=\answer{\underline{5\times 5}\times 5 = 25\times 5 = 125}$
    \end{itemize}
\end{exo}

\begin{remarque}
    Les mots \bb{carré} et \bb{cube} viennent de la géométrie. En effet
    \begin{itemize}
        \item L'aire d'un \bb{carré} de côté $c$ est $c\times c=c^2$.
        \item Le volume d'un \bb{cube} d'arête $a$ est $a\times a\times a=a^3$.
    \end{itemize}
\end{remarque}

\subsection{Carré et cube dans les priorités opératoires}
\begin{convention}
    \begin{minipage}{0.58\textwidth}
    Dans un calcul, on effectue dans l'ordre : 
    \begin{enumerate}
        \item les parenthèses
        \item les puissances (ou exposants)
        \item les multiplications et divisions
        \item les additions et soustractions
    \end{enumerate}
    \end{minipage}
    \hfill
    \vline
    \hfill
    \begin{minipage}{0.38\textwidth}
        \textcolor{Jaune}{\bb{P} : Parenthèses $()$}\\
        \textcolor{Vert}{\bb{E} : Exposants $x^2$}\\
        \textcolor{Rouge}{\bb{M} : Multiplication $\times$}\\
        \textcolor{Rouge}{\bb{D} : Division $\div$ }\\
        \textcolor{Bleu}{\bb{A} : Addition $+$}\\
        \textcolor{Bleu}{\bb{S} : Soustraction $-$}
    \end{minipage}
    \newline
    \newline
    \newline
    Il suffit de retenir l'acronyme \bb{PEMDAS}
\end{convention}


%\begin{exo}
%    Calculer en détaillant les étapes (transformer les puissances en produits,
%    puis utiliser les priorités opératoires):
%    \begin{itemize}
%        \item Exemple :
%            $2^2\times3=\underline{2\times2}\times3=\underline{4\times3}=\hl{12}$
%        \item $A = 1+5^2=$
%        \item $B = (1+6)^2=$
%        \item $C = 6\times3^2=$
%        \item $D = 1.5-1.5^2=$
%        \item $E = (2\times 3)^3=$
%    \end{itemize}
%\end{exo}


\begin{exe}
    \begin{itemize}
        \item $3+2^2=\answer{3+4=7}$  \hspace{1.5cm}mais  $(3+2)^2=\answer{5^2=25}$.
        \item $2\times5^2=\answer{2\times25=50}$  \hspace{1.5cm}mais
            $(2\times5)^2=\answer{10^2=100}$
        \item $10-2^3=\answer{10-8=2}$
        \item $4^2-3^2=\answer{16-9=7}$ \hspace{1.5cm}mais
            $(4-3)^2=\answer{1^2=1}$.
    \end{itemize}
\end{exe}

\begin{exo}
    \begin{itemize}
        \item $4+3^2=\answer{4+9 = 13}$
        \item $(4+3)^2=\answer{7^2=49}$
        \item $2\times6^2=\answer{2\times36=72}$
        \item $10-2^3=\answer{10-8=2}$
        \item $(2\times3)^3=\answer{(6)^3=216}$
    \end{itemize}
\end{exo}
